\documentclass[10pt,a4paper]{amsart}
\usepackage[margin=19mm]{geometry}
\usepackage{amsmath,amssymb,mathtools}
\usepackage[scr=rsfs,cal=euler]{mathalfa}
\usepackage{xfrac}
\usepackage[unicode,hidelinks,bookmarks=false]{hyperref}
\newcommand{\ie}{i.e.\ }
\newcommand{\cf}{cf.\ }
\newcommand{\resp}{respectively\ }
\newcommand{\Subsection}{\subsection}
\providecommand{\nobreakdash}{\nobreak\hbox{-}}
\setlength{\emergencystretch}{2em}
\multlinegap0pt
\usepackage{bm}
\usepackage{bbm}
\usepackage{dsfont}
\usepackage{xspace}
\usepackage{cancel}
\usepackage{mathtools}
\usepackage{amsthm}
\makeatletter
\@ifundefined{lemma}{\newtheorem{lemma}{Lemma}}{}
\@ifundefined{theorem}{\newtheorem{theorem}{Theorem}}{}
\makeatother
\title[Cowen-Douglas operators on quaternionic Hilbert spaces]{Cowen-Douglas operators on quaternionic Hilbert spaces}
\author{Xiaoqi Feng, Bingzhe Hou, Kui Ji}
\date{Source: arXiv:2510.19523v1 (CC BY 4.0)}
\begin{document}
\maketitle

\noindent\emph{Source and licence.} Cowen-Douglas operators on quaternionic Hilbert spaces. Version arXiv:2510.19523v1 (CC BY 4.0), distributed under the Creative Commons Attribution 4.0 International licence, as recorded in the arXiv licence metadata for this exact version and shown on the abstract page https://arxiv.org/abs/2510.19523v1. Full rights notice: licenses/arxiv-2510.19523-CC-BY-4.0.txt.\par\smallskip

\noindent\textit{Editorial scope.} Counted unit: the Theorem whose source label is \texttt{None} in the article (source0.tex lines 1481--1485, source label \texttt{None}), delivered together with the article's own complete original proof of it (source0.tex lines 1487--1842, one \texttt{proof} environment of 356 lines). No mathematics is written, repaired or re-derived: statement and proof are the article's own text line for line. Required context retained uncounted: FR (lines 1315--1317). Every definition, display and label of those lines is retained unchanged. Omitted from this package and not redistributed: 1--1314, 1318--1480, 1486--1486, 1843--2038. Nothing inside a retained excerpt is omitted or altered: every retained body is delivered exactly as the article's lines are written, so no omission receipt of any kind accompanies this package. Citations: the retained excerpts carry no citation command at all, so no bibliography entry of the article is transcribed and no reference is redistributed. Reference adaptation: none was needed and none was made; every reference inside the delivered unit resolves inside it, so no unresolved reference or citation is produced. Printed slips kept literal: no printed word, symbol, hypothesis or number of the counted statement or of its proof was altered. Layout and class: the article's own class is \texttt{sn-jnl} with packages including graphicx, multirow, amsmath, amssymb, amsfonts, amsthm, mathrsfs, appendix, xcolor, textcomp, manyfoot, booktabs, algorithm, algorithmicx; the wrapper is the lane's pinned \texttt{amsart} editorial wrapper at 10pt on A4, which restates the theorem-like environments the retained text needs and adds the article's own macro definitions that the retained text needs (none), with extra wrapper packages (bm, bbm, dsfont, xspace, cancel, mathtools). The delivered numbering is the unit's own. Assets: no retained span contains a figure environment, a graphics-inclusion command, a PDF-page inclusion, a table or a TikZ picture; no image file is copied into this package, the package declares no asset dependency and no image file is redistributed with it. Licence: version arXiv:2510.19523v1 of this article is distributed under the Creative Commons Attribution 4.0 International licence, as recorded in the arXiv licence metadata for this exact version and shown on the abstract page https://arxiv.org/abs/2510.19523v1. A full rights notice is delivered at licenses/arxiv-2510.19523-CC-BY-4.0.txt. Characters: every retained excerpt is pure ASCII, so the delivered engine, XeLaTeX with its default Latin Modern text font, reports no missing character.
    \begin{lemma}\label{FR}
 Let $T\in B_1^{s}(\Omega_q)$. Suppose that $\gamma$ and $\widetilde{\gamma}$ are two right holomorphic  frames  of $E_T$ over a domain $\Delta$. Then, there exists a complex-valued holomorphic function $f$ on $\Delta$ such that $\widetilde{\gamma}(\omega)=\gamma(\omega)f(\omega)$ for any $\omega\in\Delta$.
    \end{lemma}

\begin{theorem}
Let $T,\widetilde{T}\in B_1^{s}(\Omega_{q})$. Then,
$T$ and $\widetilde{T}$ are quaternion unitarily equivalent if and only if
$T_{\mathbb{C}}$ and $\widetilde{T}_{\mathbb{C}}$ are unitarily equivalent as bounded linear operators acting on $({\mathcal{H}_q})_\mathbb{C}$.
\end{theorem}

    \begin{proof}
    Suppose that $U$ is a quaternionic unitary operator such that
    $\widetilde{T}=UTU^*$. Notice that $(U^*)_{\mathbb{C}}=(U)_{\mathbb{C}}^*$ and $(U)_{\mathbb{C}}$ is a unitary operator. Then, it follows from $(\widetilde{T})_{\mathbb{C}}=(U)_{\mathbb{C}}(T)_{\mathbb{C}}(U)_{\mathbb{C}}^*$  that $T_{\mathbb{C}}$ and $\widetilde{T}_{\mathbb{C}}$ are unitarily equivalent as bounded linear operators acting on $({\mathcal{H}_q})_\mathbb{C}$.


    Now, suppose that there is a unitary operator $W$ acting on $({\mathcal{H}_q})_\mathbb{C}$ such that $\widetilde{T}_{\mathbb{C}}=WT_{\mathbb{C}}W^*$. Given any $\omega_0\in\Omega_{red}$. Let
    $$
    \gamma(\omega)=\alpha(\omega)+\boldsymbol{j}\beta(\omega)=
    (\alpha_{1}(\omega),\ \alpha_{2}(\omega),\ \cdots)+\boldsymbol{j}(\beta_{1}(\omega),\ \beta_{2}(\omega),\cdots)
    $$
    be a non-vanishing right holomorphic cross-section of $E_T$ over a neighborhood $\Delta$ of $\omega_0$, where $\alpha(\omega)$ and $\beta(\omega)$ are complex vectors. For any $\omega\in\Delta$, one can see that
    $$
    \boldsymbol{0}=((T-\boldsymbol{I}\omega)\gamma(\omega))_{\mathbb{C}}=(T_{\mathbb{C}}-\omega)
    \begin{pmatrix}
    	\alpha(\omega) \\
    	\beta(\omega)
    	\end{pmatrix}.
    $$
    Denote
    $$
    W
    \begin{pmatrix}
    	\alpha(\omega) \\
    	\beta(\omega)
    \end{pmatrix} \\
    =\begin{pmatrix}
    	\widetilde{\alpha}(\omega) \\
    	\widetilde{\beta}(\omega)
    \end{pmatrix}.
$$
and let
    $$
    \widetilde{\gamma}(\omega)=\widetilde{\alpha}(\omega)+\boldsymbol{j}\widetilde{\beta}(\omega).
    $$
   Since
    $$\begin{aligned}
    (\widetilde{T}-\boldsymbol{I}\omega)_{\mathbb{C}}
    \begin{pmatrix}
    \widetilde{\alpha}(\omega) \\
    \widetilde{\beta}(\omega)
    \end{pmatrix}
    = & (\widetilde{T}_{{\mathbb{C}}}-\omega)W
    \begin{pmatrix}
    \alpha(\omega) \\
    \beta(\omega)
    \end{pmatrix}\\
    = & W(T_{{\mathbb{C}}}-\omega)
    \begin{pmatrix}
    \alpha(\omega) \\
    \beta(\omega)
    \end{pmatrix}\\
    =&\boldsymbol{0}
    \end{aligned},$$
    we have
    $(\widetilde{\gamma}(\omega))_{\mathbb{C}}\in\ker{(\widetilde{T}-\boldsymbol{I}\omega)_{\mathbb{C}}}$ and consequently, $\widetilde{\gamma}(\omega)\in\ker{(\widetilde{T}-\boldsymbol{I}\omega)}$.
Notice that
\[
(\gamma(\omega)\boldsymbol{j})_{\mathbb{C}}=\begin{pmatrix}
	-\overline{\beta(\omega)} \\
	\overline{\alpha(\omega)}
\end{pmatrix} \ \ \ \text{and} \ \ \ \ (\widetilde{\gamma}(\omega)\boldsymbol{j})_{\mathbb{C}}=\begin{pmatrix}
-\overline{\widetilde{\beta}(\omega)} \\
\overline{\widetilde{\alpha}(\omega)}
\end{pmatrix}.
\]
Since
\[
({T}-\boldsymbol{I}\overline{\omega})(\gamma(\omega)\boldsymbol{j})=\gamma(\omega)\omega\boldsymbol{j}-\gamma(\omega)\boldsymbol{j}\overline{\omega}=0,
\]
we have
    $$
    \begin{aligned}
  (\widetilde{T}_{{\mathbb{C}}}-\overline{\omega})W\begin{pmatrix}
-\overline{\beta(\omega)} \\
\overline{\alpha(\omega)}
\end{pmatrix}
	= & W({T}_{\mathbb{C}}-\boldsymbol{I}\overline{\omega})\begin{pmatrix}
		-\overline{\beta(\omega)} \\
		\overline{\alpha(\omega)}
	\end{pmatrix} \\
	= & W({T}-\overline{\omega})_{\mathbb{C}}(\gamma(\omega)\boldsymbol{j})_{\mathbb{C}} \\
	= & W(({T}-\boldsymbol{I}\overline{\omega})(\gamma(\omega)\boldsymbol{j}))_{\mathbb{C}} \\
	= & \boldsymbol{0}.
\end{aligned}
$$
Then, $\left( W\begin{pmatrix}
	-\overline{\beta(\omega)} \\
	\overline{\alpha(\omega)}
\end{pmatrix} \right)_q\in\ker(\widetilde{T}-\boldsymbol{I}\overline{\omega})$. Since
\[
(\widetilde{T}-\boldsymbol{I}\overline{\omega})(\widetilde{\gamma}(\omega)\boldsymbol{j})=0,
\]
it follows from Lemma \ref{FR} that there is a complex holomorphic function $f$ on $\Delta$ such that
    $$
    W\begin{pmatrix}
    	-\overline{\beta(\omega)} \\
    	\overline{\alpha(\omega)}
    \end{pmatrix}
    =\begin{pmatrix}
    	-\overline{\widetilde{\beta}(\omega)f(\omega)} \\
    	\overline{\widetilde{\alpha}(\omega)f(\omega)}
    \end{pmatrix}.
$$
    Since $W$ is a unitary operator, one can see that
    \[
     \left\|\begin{pmatrix}
     	-\overline{\beta(\omega)} \\
     	\overline{\alpha(\omega)}
     \end{pmatrix} \right\|=\left\|W\begin{pmatrix}
     -\overline{\beta(\omega)} \\
     \overline{\alpha(\omega)}
 \end{pmatrix} \right\|=\left\|\begin{pmatrix}
 -\overline{\widetilde{\beta}(\omega)f(\omega)} \\
 \overline{\widetilde{\alpha}(\omega)f(\omega)}
\end{pmatrix} \right\|=\left\|\begin{pmatrix}
-\overline{\beta(\omega)} \\
\overline{\alpha(\omega)}
\end{pmatrix} \right\|\cdot|f(\omega)|,
    \]
  and consequently for every $\omega\in\Delta$,
    \[
    |f(\omega)|=1.
    \]
   Then, there exists $\theta_0\in[0,\ 2\pi)$ such that
    $f(\omega)=\mathrm{e}^{\boldsymbol{i}\theta_0}$ for every $\omega\in\Delta$. Define the linear operator $V:(\mathcal{H}_q)_{\mathbb{C}}\rightarrow(\mathcal{H}_q)_{\mathbb{C}}$ by
    $$
    V
    \begin{pmatrix}
    \alpha^{(l)}(\omega_0) \\
    \beta^{(l)}(\omega_0) \\
    \end{pmatrix}
    =\begin{pmatrix}
    \widetilde{\alpha}^{(l)}(\omega_0)\mathrm{e}^{-\boldsymbol{i}\theta_0/2} \\
    \widetilde{\beta}^{(l)}(\omega_0)\mathrm{e}^{-\boldsymbol{i}\theta_0/2} \\
    \end{pmatrix},\
    V
    \begin{pmatrix}
    -\overline{\beta^{(l)}(\omega_0)} \\
    \overline{\alpha^{(l)}(\omega_0)} \\
    \end{pmatrix}
    =\begin{pmatrix}
    -\overline{\widetilde{\beta}^{(l)}(\omega_0)}\mathrm{e}^{\boldsymbol{i}\theta_0/2} \\
    \overline{\widetilde{\alpha}^{(l)}(\omega_0)}\mathrm{e}^{\boldsymbol{i}\theta_0/2} \\
    \end{pmatrix}
    $$
    for any $l\in\mathbb{N}$. It is not difficult to see that $VT_{\mathbb{C}}=\widetilde{T}_{\mathbb{C}}V$.

  For every $k,\ l\in\mathbb{N}$, we have
     \begin{equation*}
 	\begin{aligned}
 		\langle V
 		\begin{pmatrix}
 			\alpha^{(l)}(\omega_0) \\
 			\beta^{(l)}(\omega_0)
 		\end{pmatrix},\
 		V
 		\begin{pmatrix}
 			\alpha^{(k)}(\omega_0) \\
 			\beta^{(k)}(\omega_0)
 		\end{pmatrix}
 		\rangle
 		= & \langle
 		\begin{pmatrix}
 			\widetilde{\alpha}^{(l)}(\omega_0)\mathrm{e}^{-\boldsymbol{i}\theta_0/2} \\
 			\widetilde{\beta}^{(l)}(\omega_0)\mathrm{e}^{-\boldsymbol{i}\theta_0/2} \\
 		\end{pmatrix},\
 		\begin{pmatrix}
 			\widetilde{\alpha}^{(k)}(\omega_0)\mathrm{e}^{-\boldsymbol{i}\theta_0/2} \\
 			\widetilde{\beta}^{(k)}(\omega_0)\mathrm{e}^{-\boldsymbol{i}\theta_0/2} \\
 		\end{pmatrix}
 		\rangle \\
 		= &
 		\langle
 		\begin{pmatrix}
 			\widetilde{\alpha}^{(l)}(\omega_0) \\
 			\widetilde{\beta}^{(l)}(\omega_0)\\
 		\end{pmatrix},\
 		\begin{pmatrix}
 			\widetilde{\alpha}^{(k)}(\omega_0) \\
 			\widetilde{\beta}^{(k)}(\omega_0) \\
 		\end{pmatrix}
 		\rangle \\
 		= & \langle
 		W\begin{pmatrix}
 			\alpha^{(l)}(\omega_0) \\
 			\beta^{(l)}(\omega_0)
 		\end{pmatrix},\
 		W\begin{pmatrix}
 			\alpha^{(k)}(\omega_0) \\
 			\beta^{(k)}(\omega_0)
 		\end{pmatrix}
 		\rangle \\
 		= & \langle
 		\begin{pmatrix}
 			\alpha^{(l)}(\omega_0) \\
 			\beta^{(l)}(\omega_0)
 		\end{pmatrix},\
 		\begin{pmatrix}
 			\alpha^{(k)}(\omega_0) \\
 			\beta^{(k)}(\omega_0)
 		\end{pmatrix}
 		\rangle,
 	\end{aligned}
 \end{equation*}

    \begin{equation*}
    \begin{aligned}
     \langle V
    \begin{pmatrix}
    -\overline{\beta^{(l)}(\omega_0)} \\
    \overline{\alpha^{(l)}(\omega_0)}
    \end{pmatrix},\
    V
   \begin{pmatrix}
   	-\overline{\beta^{(k)}(\omega_0)} \\
   	\overline{\alpha^{(k)}(\omega_0)}
   \end{pmatrix}
    \rangle
    = &      \langle
    \begin{pmatrix}
    	-\overline{\widetilde{\beta}^{(l)}(\omega_0)}\mathrm{e}^{\boldsymbol{i}\theta_0/2} \\
    	\overline{\widetilde{\alpha}^{(l)}(\omega_0)}\mathrm{e}^{\boldsymbol{i}\theta_0/2}
    \end{pmatrix},\
      \begin{pmatrix}
     	-\overline{\widetilde{\beta}^{(k)}(\omega_0)}\mathrm{e}^{\boldsymbol{i}\theta_0/2} \\
     	\overline{\widetilde{\alpha}^{(k)}(\omega_0)}\mathrm{e}^{\boldsymbol{i}\theta_0/2}
     \end{pmatrix}
    \rangle \\
    = &
    \langle
    \begin{pmatrix}
    	-\overline{\widetilde{\beta}^{(l)}(\omega_0)}\mathrm{e}^{\boldsymbol{i}\theta_0} \\
    	\overline{\widetilde{\alpha}^{(l)}(\omega_0)}\mathrm{e}^{\boldsymbol{i}\theta_0}
    \end{pmatrix},\
    \begin{pmatrix}
    	-\overline{\widetilde{\beta}^{(k)}(\omega_0)}\mathrm{e}^{\boldsymbol{i}\theta_0} \\
    	\overline{\widetilde{\alpha}^{(k)}(\omega_0)}\mathrm{e}^{\boldsymbol{i}\theta_0}
    \end{pmatrix}
    \rangle \\
    = & \langle W
    \begin{pmatrix}
    	-\overline{\beta^{(l)}(\omega_0)} \\
    	\overline{\alpha^{(l)}(\omega_0)}
    \end{pmatrix},\
    W
    \begin{pmatrix}
    	-\overline{\beta^{(k)}(\omega_0)} \\
    	\overline{\alpha^{(k)}(\omega_0)}
    \end{pmatrix}
    \rangle \\
     = & \langle
     \begin{pmatrix}
     	\alpha^{(l)}(\omega_0) \\
     	\beta^{(l)}(\omega_0)
     \end{pmatrix},\
       \begin{pmatrix}
     	\alpha^{(k)}(\omega_0) \\
     	\beta^{(k)}(\omega_0)
     \end{pmatrix}
     \rangle,
    \end{aligned}
    \end{equation*}
and
     \begin{equation*}
	\begin{aligned}
		\langle V
		\begin{pmatrix}
			\alpha^{(l)}(\omega_0) \\
			\beta^{(l)}(\omega_0)
		\end{pmatrix},\
		 V
		\begin{pmatrix}
			-\overline{\beta^{(k)}(\omega_0)} \\
			\overline{\alpha^{(k)}(\omega_0)}
		\end{pmatrix}
		\rangle
		= & \langle
		\begin{pmatrix}
			\widetilde{\alpha}^{(l)}(\omega_0)\mathrm{e}^{-\boldsymbol{i}\theta_0/2} \\
			\widetilde{\beta}^{(l)}(\omega_0)\mathrm{e}^{-\boldsymbol{i}\theta_0/2} \\
		\end{pmatrix},\
		\begin{pmatrix}
			-\overline{\widetilde{\beta}^{(k)}(\omega_0)}\mathrm{e}^{\boldsymbol{i}\theta_0/2} \\
			\overline{\widetilde{\alpha}^{(k)}(\omega_0)}\mathrm{e}^{\boldsymbol{i}\theta_0/2}
		\end{pmatrix}
		\rangle \\
		= &
		\langle
		\begin{pmatrix}
			\widetilde{\alpha}^{(l)}(\omega_0) \\
			\widetilde{\beta}^{(l)}(\omega_0)\\
		\end{pmatrix},\
		\begin{pmatrix}
			-\overline{\widetilde{\beta}^{(k)}(\omega_0)}\mathrm{e}^{\boldsymbol{i}\theta_0} \\
			\overline{\widetilde{\alpha}^{(k)}(\omega_0)}\mathrm{e}^{\boldsymbol{i}\theta_0}
		\end{pmatrix}
		\rangle \\
		= & \langle
		W\begin{pmatrix}
			\alpha^{(l)}(\omega_0) \\
			\beta^{(l)}(\omega_0)
		\end{pmatrix},\
		W\begin{pmatrix}
			-\overline{\beta^{(k)}(\omega_0)} \\
			\overline{\alpha^{(k)}(\omega_0)}
		\end{pmatrix}
		\rangle \\
		= & \langle
		\begin{pmatrix}
			\alpha^{(l)}(\omega_0) \\
			\beta^{(l)}(\omega_0)
		\end{pmatrix},\
		\begin{pmatrix}
			-\overline{\beta^{(k)}(\omega_0)} \\
			\overline{\alpha^{(k)}(\omega_0)}
		\end{pmatrix}
		\rangle,
	\end{aligned}
\end{equation*}

   Then, $V$ is an isometry. Following from
   \[
   {\bigvee}_{\mathbb{C}}\{\widetilde{\gamma}^{(i)}(\omega_0); i\in\mathbb{N}\ \}=\mathcal{H}_q,
   \]
   $V$ is surjective ahd hence $V$ is a unitary operator satisfying $\widetilde{T}_{\mathbb{C}}=VT_{\mathbb{C}}V^*$.

Define the linear operator $U:\mathcal{H}_q\rightarrow\mathcal{H}_q$ by for any $\omega\in\Delta$,
\[
U(\gamma(\omega))=\widetilde{\gamma}(\omega)\cdot \mathrm{e}^{-\boldsymbol{i}\theta_0/2}.
\]
Notice that for any $l\in\mathbb{N}$,
\[
U_{\mathbb{C}}
\begin{pmatrix}
	\alpha^{(l)}(\omega_0) \\
	\beta^{(l)}(\omega_0) \\
\end{pmatrix}
=(U\gamma^{(l)}(\omega_0))_{\mathbb{C}}=(\widetilde{\gamma}^{(l)}(\omega_0)\mathrm{e}^{-\boldsymbol{i}\theta_0/2})_{\mathbb{C}}=\begin{pmatrix}
	\widetilde{\alpha}^{(l)}(\omega_0)\mathrm{e}^{-\boldsymbol{i}\theta_0/2} \\
	\widetilde{\beta}^{(l)}(\omega_0)\mathrm{e}^{-\boldsymbol{i}\theta_0/2} \\
\end{pmatrix},
\]
and
\[
U_{\mathbb{C}}
\begin{pmatrix}
	-\overline{\beta^{(l)}(\omega_0)} \\
	\overline{\alpha^{(l)}(\omega_0)} \\
\end{pmatrix}
=(U\gamma^{(l)}(\omega_0)\boldsymbol{j})_{\mathbb{C}}=(\widetilde{\gamma}^{(l)}(\omega_0)\mathrm{e}^{-\boldsymbol{i}\theta_0/2}\boldsymbol{j})_{\mathbb{C}}=\begin{pmatrix}
	-\overline{\widetilde{\beta}^{(l)}(\omega_0)}\mathrm{e}^{\boldsymbol{i}\theta_0/2} \\
	\overline{\widetilde{\alpha}^{(l)}(\omega_0)}\mathrm{e}^{\boldsymbol{i}\theta_0/2} \\
\end{pmatrix}.
\]
Then, $U_{\mathbb{C}}=V$ and hence $U$ is a quaternionic unitary operator with $\widetilde{T}=UTU^*$. This finishes the proof.
\end{proof}

\end{document}
